\chapter{Generell informasjon}

Masteroppgaven er siste steg mot en mastergrad i Multimedia og pedagogisk teknologi ved Universitetet i Agder, Fakultet for teknologi og realfag.

En fullført oppgave utgjør 30 studiepoeng og gjennomføres i fjerde og siste semester.

\section{MM-500 Emnebeskrivelse}

For å kvalifisere til å skrive masteroppgave i MM-500 må du ha bestått eksamener tilsvarende 70 studiepoeng innen utdanningsplanen ved starten av semesteret hvor oppgaven skal skrives. Emner som er grunnleggende for oppgaven må være bestått. Tema for oppgaven må godkjennes av Universitetet i Agder.

\subsection{Emneinnhold}

Masteroppgaven er et selvstendig arbeid med tema relevant for programmets profil. Arbeidet skal inneholde forskning og ny kunnskap eller nye metoder. Grunnleggende oppgaver innen planlegging, programmering eller utvikling bør unngås.

Prosjektet inkluderer:
\begin{itemize}
\item Rapport som beskriver forskningsstatus, problemstilling, tilnærming, resultater, evaluering og videre perspektiver
\item Pressemelding som presenterer oppgaven og hovedresultatene
\item Plenums­presentasjon
\end{itemize}
